\chapter{Cascaded Pictorial Structures}\label{sec:CPS}

Pictorial structure models~\cite{fischler1973ps} are a popular method for human body pose estimation~\cite{felz05,fergus2005sparse,devacrf,ferrari08,andriluka09}.
The model is a Conditional Random Field over pose variables that characterizes 
local appearance properties of parts and geometric part-part interactions.   
The search over the joint pose space is linear time in the number of parts when 
the part-part dependencies form a tree.  However,  the individual part 
state-spaces are too large (typically hundreds of thousands of states) to allow 
complex appearance models be evaluated densely.   Most appearance models are 
therefore simple linear filters on edges, color and 
location~\cite{felz05,devacrf,ferrari08,andriluka09}. Similarly, because of 
quadratic state-space complexity, part-part relationships are typically 
restricted to be image-independent deformation costs that allow for convolution 
or distance transform tricks to speed up inference~\cite{felz05}. A common 
problem in such models is poor localization of parts that have weak appearance 
cues or are easily confused with background clutter (accuracy for lower arms in 
human figures is almost half of that for torso or head~\cite{andriluka09}).   
Localizing these elusive parts requires richer models of individual part shape 
and joint part-part appearance, including contour continuation and segmentation 
cues, which are prohibitive to compute densely.

In order to enable richer appearance models, we propose to learn a cascade of 
pictorial structures (\CPS) of increasing pose resolution which 
progressively filter the pose state space.  Conceptually, the idea is similar 
to the work on cascades for face detection~\cite{geman2001,viola02}, but the 
key difference is the use of structured models. Each level of the cascade at a 
given spatial/angular resolution refines the set of candidates from the 
previous level and then runs inference to determine which poses to filter out.  
For each part, the model selects poses with the largest max-marginal scores, 
subject to a computational budget.  Unlike conventional pruning heuristics, 
where the possible part locations are identified using the output of a 
detector, models in our cascade use inference in simpler structured models to 
identify what to prune, taking into account global pose in filtering decisions.
%\begin{figure}[t]
%\begin{center}
%\includegraphics[width=0.8\textwidth]{figs/empty.jpg}
%\end{center}
%\caption{\label{fig:overview} Overview: A discriminative coarse-to-fine cascade 
%of pictorial structures filters the pose space so that expressive and 
%computationally expensive cues can be used in the final pictorial structure.  
%Shown are 5 levels of our coarse-to-fine cascade for the right upper and lower 
%arm parts.  Green vectors represent position and angle of unpruned states, the 
%downsampled images correspond to the dimensions of the resepective state space, 
%and the white rectangles represent classification using our final model.}
%\end{figure}
As a result, at the final level the \CPS{} model has to deal with a much 
smaller hypotheses set which allows us to use a rich combination of features.  
In addition to the traditional part detectors and geometric features, we are 
able to incorporate object boundary continuity and smoothness, as well as shape 
features. The former features represent mid-level and bottom-up cues, while the 
latter capture shape information, which is complementary to the traditional 
HoG-based part models.  The approach is illustrated in the overview 
Figure~\ref{fig:overview}. We apply the presented \CPS model combined 
with the richer set of features on the Buffy and PASCAL stickmen benchmark, 
improving the state-of-the-art on arm localization. 

\label{subsec:our_ps}
We choose instead to model part configurations as a general log-linear Conditional Random Field over pairwise and unary terms:
\begin{align}
p( y | x) &\propto \exp \Biggl[\sum_{ij} \w_{ij} \cdot \f_{ij}(y_i,y_j,x) +
\sum_i \w_i \cdot \f_i(y_i,x) \Biggl] = e^{\w \cdot \f(y,x)}.
\end{align}
The parameters of our model are the pairwise and unary weight vectors 
$\theta_{ij}$ and $\theta_i$ corresponding to the pairwise and unary feature 
vectors $\phi_{ij}(y_i,y_j,x)$ and $\phi_i(\l_i,x)$.  For brevity, we stack all 
the parameters and features into vectors using notation $\theta^T \phi(ys,x)$.  
The key differences with the classical PS model are that (1) our pairwise costs 
allow data-dependent terms, and (2) we do not constrain our parameters to fit 
any parametric distribution such as a Gaussian.  For example, we can express 
the pairwise features used in the classical model as $y_{i} \cdot y_{i} \text{, 
} y_{j}\cdot y_{j} \text{ and }  y_{i}\cdot y_{j}$ without requiring that their 
corresponding weights can be combined into a positive semi-definite covariance 
matrix.

In this general form, inference can not be performed efficiently with distance transforms or convolution, and we rely on standard $O(|\mathcal{Y}_i|^2)$ dynamic programming techniques to compute the MAP assignment or part posteriors.  Many highly-effective pairwise features one might design would be intractable to compute in this manner for a reasonably-sized state space---for example an $100\times100$ image with a part angle discretization of $24$ bins yields $|\mathcal{Y}_i|^2 = 57.6$ billion part-part hypotheses.

In the next section, we describe how we circumvent this issue via a cascade of models which aggressively prune the state space at each stage typically without discarding the correct sequence.  After the state space is pruned, we are left with a small enough number of states to be able to incorporate powerful data-dependent pairwise and unary features into our model.


\subsection*{Structured Prediction Cascades} \label{cascades}
The recently introduced Structured Prediction Cascade framework~\cite{cascades} provides a principled way to prune the state space of a structured prediction problem via a sequence of increasingly complex models.
There are many possible ways of defining a sequence of increasingly complex models.  In~\cite{cascades} the authors introduce higher-order cliques into their models in successive stages (first unary, then pairwise, ternary, etc.).  Another option is to start with simple but computationally efficient features, and add more complex features downstream as the number of states decreases.  Yet another option is to geometrically coarsen the original state space and successively prune and refine.
We use a coarse-to-fine state space approach with simple features until we are at a reasonably fine enough state space resolution and left with few enough states that we can introduce more complex features.  We start with a severely coarsened state space and use standard pictorial structures unary detector scores and geometric features to perform quick exhaustive inference on the coarse state space.  

%The simplest analogy of the framework is to the popular cascade classifier introduced by Viola and Jones for binary prediction problems~\cite{Viola2002}.
More specifically, each level of the cascade uses inference to identify which states to prune away and the next level refines the spatial/angular resolution on the unpruned states.
The key ingredient to the cascade framework is that states are pruned using {\em max-marginal} scores, computed using dynamic programming techniques.  For brevity of notation, define the score of a joint part state $ys$ as $\theta_x(ys)$ and the max-marginal score of a part state as follows:
\begin{align}
\theta_x(ys) &= \theta^T \phi(ys,x)= \sum_{ij}  \theta_{ij}^T \phi_{ij}(y_i,y_j,x) +
\sum_i \theta_i^T\phi_i(y_i,x) \\
\mmi  &= \max_{y' \in ys} \;\;\; \{ \theta_x(y') \; : \; y_i'=y_i \}
\end{align}

In words, the max-marginal for location/angle $y_i$ is the score of the best sequence which constrains $y_i = y_i$.  In a pictorial structure model, this corresponds to fixing limb $i$ at location $y_i$, and determining the highest scoring configuration of other part locations and angles under this constraint.  A part could have weak individual image evidence of being at location $y_i$ but still have a high max-marginal score if the rest of the model believes this is a likely location.  Similarly, we denote the MAP assignment score as
$\theta_x^\star = \max_{ys \in ys} \theta_x(ys)$, the unconstrained best configuration of all parts.

When learning a cascade, we have two competing objectives that we must trade off,
accuracy and efficiency: we want to minimize the number of errors incurred by each
level of the cascade  and maximize the number of filtered max marginals.  A natural strategy is to prune away the lowest ranked states based on max-marginal scores. Instead, \cite{cascades} prune the states whose max-marginal score is lower than an data-specific threshold $t_x$: $l_i$ is pruned if $\mmi < t_x$.  This threshold is defined as a convex combination of the MAP assignment score and the mean max-marginal score, meant to approximate a percentile threshold: $$t_x(\theta,\alpha) = \alpha \theta^\star_x + (1-\alpha)\frac{1}{M}\sum_{i=1}^M\frac{1}{|\mathcal{Y}_i|}\sum_{y_i \in \mathcal{Y}_i} \mmi,$$ where $\alpha\in[0,1]$ is a parameter to be chosen that determines how aggressively to prune. When $\alpha = 1$, only the best state is kept, which is equivalent to finding the MAP assignment.  When $\alpha = 0$ approximately half of the states are pruned (if the median of max-marginals is equal to the mean) .  The advantage of using  $t_x(\theta,\alpha)$ is that it is convex in $\theta$, and leads to a convex formulation for parameter estimation that trades off the proportion of incorrectly pruned states with the proportion of unpruned states.
Note that  $\alpha$ controls efficiency, so we focus on learning the parameters $\theta$ that 
minimize the number of errors for a given filtering level $\alpha$.   The learning formulation uses a simple fact about max-marginals and the definition of $t_x(\theta,\alpha)$
to get a handle on errors of the cascade:
if $\theta_x(ys) > t_x(\theta,\alpha)$, then for all i, $\mmi >  t_x(\theta,\alpha)$, so no part state of $ys$ is pruned.   Given an example $(x,ys)$,
this condition $\theta_x(ys) > t_x(\theta,\alpha)$ is sufficient to ensure that no correct part is pruned.
\begin{figure}[t]
\begin{center}
\includegraphics[width=0.75\textwidth]{figs/empty.jpg}
\caption[SHORT TITLE]{Upper right: Detector-based pruning by thresholding (for 
the lower right arm) yields many hypotheses far way from the true one. Lower 
row: The CPS, however, exploits global information to perform better pruning.}
\label{fig:cascade_pruning}
\end{center}
\end{figure}
%\end{floatingfigure}

To learn one level of the structured cascade model $\theta$ for a fixed $\alpha$, we try to minimize 
the number of correct states that are pruned on training data by solving the
following convex margin optimization problem given $N$ training examples $(x^n,ys^n)$:
\begin{equation}
  \label{eq:convex_opt}
 \min_{\theta} \;\;\;\ \frac{\lambda}{2}||\theta||^2 + \frac{1}{N}\sum_{n=1}^N H(\theta;x^n,ys^n),
\end{equation}
where $H$ is a hinge upper bound 
$H(\theta;x,ys) = \max\{0, 1 + t_x(\theta,\alpha) - \theta_x(ys)\}$.
The upper-bound $H$ is a hinge loss measuring the margin between the
filter threshold $t_{x^n}(\theta,\alpha)$ and the score of the truth
$\theta^T\phi(ys^n,x^n)$; the loss is zero if the truth scores above
the threshold by margin $1$.  We solve \eqref{eq:convex_opt} using stochastic
sub-gradient descent. Given an example $(x,ys)$, we apply the following
update if $H(\theta;x,ys)$ (and the sub-gradient) is non-zero:
\begin{equation*}
    \label{eq:spf_update}
    \theta'  \leftarrow \theta + \eta \left(-\lambda \theta + \phi(ys,x) - \alpha \phi(ys^\star,x)  - (1-\alpha) \frac{1}{M} \sum_{i} \frac{1}{|\mathcal{Y}_i|}\sum_{y_i \in \mathcal{Y}_i}  \phi(ys^\star(ys_i),x)\right).
\end{equation*}
Above, $\eta$ is a learning rate parameter, $ys^\star = \argmax_{ys'}
\theta_x(ys')$ is the highest scoring assignment and $ys^\star(ys_i) = \argmax_{ys': ys'_i = l_i} \theta_x(ys')$
are highest scoring assignments constrained to $y_i$ for part $i$.
The key distinguishing feature of this update as compared to
structured perceptron is that it subtracts features included in all
max-marginal assignments $ys^\star(ys_i)$\footnote{Note that because \eqref{eq:convex_opt} is $\lambda$-strongly convex,
if we chose $\eta_t = 1/(\lambda t)$ and add a projection step to keep 
$\theta$ in a closed set, the update would correspond to the Pegasos
update with convergence guarantees of $\tilde{O}(1/\epsilon)$ iterations for
$\epsilon$-accurate solutions~\cite{shalev-shwartz07pegasos}.  In our experiments, we found the projection step made no difference and used only 2 passes over the data, with $\eta$ fixed.}.
%, as we discuss in Section~\ref{implementation}.

The stages of the cascade are learned sequentially, from coarse to fine, and each has a different $\theta$ and $\mathcal{Y}_i$ for each part, as well as  $\alpha$.  The states of the next level are simply refined versions of the states that have not been pruned. We describe 
the refinement structure of the cascade in Section~\ref{implementation}.
In the end of a coarse-to-fine cascade we are left with a small, sparse set of states that typically contains the groundtruth states or states relatively close to them---in practice we are left with around 500 states per part, and 95\% of the time we retain a state the is close enough to be considered a match (see Table~\ref{tab:pruning}). At this point we have the freedom to add a variety of complex unary and pairwise part interaction features involving geometry, appearance, and compatibility with perceptual grouping principles which we describe in Section~\ref{sec:features}.

\mypar{Why not just detector-based pruning?} A naive approach used in a variety 
of applications is to simply subsample states by thresholding outputs of part 
or sparse feature detectors, possibly combined with non-max suppression.  Our 
approach, based on pruning on max-marginal values in a first-order model, is 
more sophisticated: for articulated parts-based models, strong evidence from 
other parts can keep a part which has weak individual evidence, and would be 
pruned using only detection scores.  The failure of prefiltering  part 
locations in human pose estimation is also noted by~\cite{andriluka09}, and 
serves as the primary justification for their use of the dense classical PS. 
This is illustrated in Figure~\ref{fig:cascade_pruning} on an example image 
from~\cite{ferrari08}.


